% Appendix A - two-page policy brief for the Director (non-technical: no code, no jargon)
\section{Policy Brief for the Director, Disease Control \& Preparedness Unit}
\begin{center}\large\textbf{Seeing outbreaks sooner and sending supplies where they save most lives}\\[2pt]
\normalsize EpiScope prototype -- summary for decision-makers\end{center}

\noindent\fbox{\parbox{0.97\textwidth}{\textbf{Important:} this prototype was built and tested on realistic
\emph{simulated} patient records for all 64 districts (2021--2025), because real patient data cannot be used
for a student project. The findings below show what the system can do; the numbers must be re-checked on the
Ministry's own data before any operational decision.}}

\subsection*{1. The current situation}
Information about dengue, diarrhoea and cholera, influenza-like illness and malaria reaches the Ministry late,
in many formats, and is mostly analysed after the fact. As a result, saline, test kits and teams often reach a
district after its hospitals are already full. EpiScope brings the routine reports from hospitals and health
complexes, the weather record, population figures and district stock levels together every day and turns them
into four things officials can act on: \textbf{early warnings}, \textbf{forecasts for the next four weeks},
\textbf{a recommended distribution of supplies and response teams}, and \textbf{a calendar for the year ahead}.

\subsection*{2. Top three recommendations}
\begin{enumerate}
\item \textbf{Adopt a weekly ``warning review'' based on the ranked warning list.} In our test the warning system
recognised \SensLocal{} of the simulated local outbreaks, typically about \MedDelayLocal{} days after they began,
while producing roughly \FAperDY{} unnecessary warnings per district per year across the four diseases. It also
ignored a hospital that suddenly sent three weeks of late reports at once -- a situation that fools simpler systems.
\emph{Expected impact:} response teams can be sent days to weeks earlier than with monthly spreadsheet reviews,
with a manageable number of false alarms. Each warning should be confirmed by a phone call to the district
before a team is sent.
\item \textbf{Distribute supplies with the recommended plan, guaranteeing every district a small minimum
share} (at least one tenth of its share by population of every shipment). With a budget of 60\% of what would be needed to fill every district's shortfall, the
recommended plan fills \OptGapCov{}\% of the districts' combined shortfall, against \BaseGapCov{}\% if supplies
are shared by population alone, and guarantees every district at least a share of each shipment. This guarantee costs a little efficiency
(about \RecEffLoss{}\% more unmet need than a plan that ignores fairness) but means no district is left with
nothing. \emph{Expected impact:} more saline and test kits where patients will actually be, without abandoning
low-risk districts.
\item \textbf{Follow the preparedness calendar: start dengue prevention between mid-March and late April, and have
supplies in place by May.} In every division dengue cases begin to rise in May or June and fall away by late
November; the peak comes first in Sylhet and Mymensingh (July) and last in Rajshahi (late October). Diarrhoea
rises from March and peaks between April and June. The calendar tells each division when to begin larval control (about eight weeks before the rise),
when supplies must be in place (four weeks before) and when to add hospital staff (two weeks before the peak).
\emph{Expected impact:} the heaviest weeks of the season are met with stock already in place.
\end{enumerate}

\subsection*{3. How to read the dashboard}
\begin{itemize}
\item \textbf{Map:} darker blue = more cases per 100{,}000 people in the chosen weeks. A red triangle marks a
district with an active warning; an orange diamond means ``check the reporting first'' (for example, a new
hospital has just joined the system). Striped areas include small numbers that are hidden for privacy.
\item \textbf{Weekly chart:} the solid line is what has been reported; the dashed end is our estimate for the
weeks still being reported; the shaded fan is the forecast -- the darkest band is where the true number is
most likely to fall, the lightest covers 9 out of 10 outcomes. The red dotted line is the warning threshold.
\item \textbf{Warning list:} districts at the top need attention first. ``Days in alert'' shows how long the
warning has lasted.
\item \textbf{Supplies chart and sliders:} compare expected need (orange), current stock (black diamond) and the
recommended shipment (blue). The sliders let you ask ``what if the budget were smaller?'' or ``what if we
guarantee every district a larger minimum share?'' and see the effect immediately.
\item Every panel shows the date of its latest data. A district officer sees their own district in full and the
rest of the country only in summary.
\end{itemize}

\subsection*{4. Confidence and limitations}
\begin{itemize}
\item Forecasts one week ahead were much closer to what happened than simply assuming ``same as this week last
year'': their error score was only \GlmRelWISOne{} times that of the simple method, and the stated ranges
contained the real number about as often as promised (8 times in 10). Forecasts become less certain further ahead, and the
fans in the chart show this honestly. Four-week forecasts should guide preparation, not exact orders.
\item Slow-building outbreaks are harder to spot than sudden ones; unusual weather (floods, an early monsoon)
also causes genuine rises that the system will flag and that officials must interpret.
\item The results depend on hospitals reporting regularly. The system detects facilities that stop reporting,
but it cannot see patients who never reach a facility.
\item Supply needs are converted from cases using standard assumptions (for example, six bags of saline per
admitted dengue patient); these should be confirmed with clinicians and stores staff.
\item Some public inputs (weather, population) must be refreshed from the official sources when the system is
installed at the Ministry.
\end{itemize}

\subsection*{5. Data-protection assurance}
\begin{itemize}
\item Names and national ID numbers never leave the secure first storage area: they are replaced by an
irreversible code made with a secret key that is kept outside the software.
\item Exact ages are reduced to age groups. The dashboard and every export show only totals per district and
week; any count between 1 and 4 is shown as ``fewer than 5'', and a second number is hidden when totals would
otherwise reveal it.
\item Access is by personal key, each user sees only what their role requires, and every access is logged.
\item Any record-level file shared with researchers is checked so that each combination of place, time, age group
and sex covers at least five people; for public releases we recommend adding a small amount of statistical noise
(differential privacy).
\end{itemize}
